
\documentclass{article}
\usepackage{colortbl}
\usepackage{makecell}
\usepackage{multirow}
\usepackage{supertabular}

\begin{document}

\newcounter{utterance}

\centering \large Interaction Transcript for game `dond', experiment `coop', episode 19 with Qwen3.6{-}35B{-}A3B{-}FP8{-}without{-}reasoning.
\vspace{24pt}

{ \footnotesize  \setcounter{utterance}{1}
\setlength{\tabcolsep}{0pt}
\begin{supertabular}{c@{$\;$}|p{.15\linewidth}@{}p{.15\linewidth}p{.15\linewidth}p{.15\linewidth}p{.15\linewidth}p{.15\linewidth}}
    \# & $\;$A & \multicolumn{4}{c}{Game Master} & $\;\:$B\\
    \hline

    \theutterance \stepcounter{utterance}  
    & & \multicolumn{4}{p{0.6\linewidth}}{
        \cellcolor[rgb]{0.9,0.9,0.9}{
            \makecell[{{p{\linewidth}}}]{
                \texttt{\tiny{[P1$\langle$GM]}}
                \texttt{Вы участвуете в игре кооперативных переговоров, в которой должны договориться с другим игроком о разделении набора предметов между вами.} \\
\\ 
\texttt{Правила:} \\
\texttt{(a) Вы и другой игрок получаете набор предметов. Каждому из вас также дается секретная функция стоимости, представляющая, как вы оцениваете каждый тип объекта.} \\
\texttt{(b) Вы обмениваетесь сообщениями с другим игроком, чтобы согласиться, кто какие предметы получает. Каждый может отправить максимум 5 сообщений или закончить игру раньше, сделав секретное предложение в любой момент.} \\
\texttt{(c) Вас просят отправить секретное предложение, указывающее предметы, которые вы хотите, в следующем формате в квадратных скобках: "{[}Предложение: <число> <названия предмета>, <число> <названия предмета>, <...>{]}"} \\
\texttt{(d) Если ваши предложения дополняют друг друга, то есть есть достаточно предметов для выполнения обоих предложений, каждый игрок получает баллы на основе суммы стоимостей предметов, которые он получил. В противном случае вы оба получаете нулевой балл.} \\
\texttt{(e) Ваша цель {-} максимизировать сумму вашего балла и балла другого игрока.} \\
\\ 
\texttt{Давайте начнем.} \\
\\ 
\texttt{Набор доступных предметов:} \\
\\ 
\texttt{2 лампы, 3 цветы, 1 банан} \\
\\ 
\texttt{Ваша секретная функция стоимости:} \\
\\ 
\texttt{лампа: 4, цветок: 0, банан: 2} \\
\\ 
\texttt{ВАЖНО: Ваши сообщения, если только это не секретное предложение, передаются прямо другому игроку, поэтому не включайте никакие ответы на правила или текст, объявляющий ваше сообщение. Чтобы сделать секретное предложение, используйте указанный формат. Не используйте квадратные скобки при общении с другим игроком, иначе это будет интерпретировано как ваше секретное предложение.} \\
            }
        }
    }
    & \\ \\

    \theutterance \stepcounter{utterance}  
    & \multicolumn{4}{p{0.6\linewidth}}{
        \cellcolor[rgb]{0.9,0.9,0.9}{
            \makecell[{{p{\linewidth}}}]{
                \texttt{\tiny{[P1$\rangle$GM]}}
                \texttt{Привет! Давай попробуем договориться.} \\
\\ 
\texttt{У меня есть следующие предпочтения:} \\
\texttt{{-} Лампа: 4 балла} \\
\texttt{{-} Банан: 2 балла} \\
\texttt{{-} Цветок: 0 баллов} \\
\\ 
\texttt{Поскольку цветы для меня ничего не стоят, я готов отдать их тебе, если ты возьмешь на себя лампы или банан.} \\
\\ 
\texttt{Предлагаю такой вариант:} \\
\texttt{{-} Я беру 1 лампу и 1 банан.} \\
\texttt{{-} Ты берешь 1 лампу, 3 цветка и 0 бананов (или 1 банан, если тебе он нужнее).} \\
\\ 
\texttt{Но чтобы максимизировать общую сумму, давай посмотрим на общую ценность:} \\
\texttt{{-} 2 лампы = 8 баллов (для меня)} \\
\texttt{{-} 3 цветка = 0 баллов (для меня)} \\
\texttt{{-} 1 банан = 2 балла (для меня)} \\
\texttt{Общая ценность для меня: 10 баллов.} \\
\\ 
\texttt{Если ты оцениваешь предметы иначе, скажи мне свои предпочтения. Например, если тебе важны цветы, я с радостью отдам их тебе.} \\
\\ 
\texttt{Жду твоего предложения или контрпредложения.} \\
            }
        }
    }
    & & \\ \\

    \theutterance \stepcounter{utterance}  
    & & & \multicolumn{2}{p{0.3\linewidth}}{
        \cellcolor[rgb]{0.9,0.9,0.9}{
            \makecell[{{p{\linewidth}}}]{
                \texttt{\tiny{[GM$|$GM]}}
                \texttt{continue} \\
            }
        }
    }
    & & \\ \\

    \theutterance \stepcounter{utterance}  
    & & & \multicolumn{2}{p{0.3\linewidth}}{
        \cellcolor[rgb]{0.9,0.9,0.9}{
            \makecell[{{p{\linewidth}}}]{
                \texttt{\tiny{[GM$|$GM]}}
                \texttt{Привет! Давай попробуем договориться.} \\
\\ 
\texttt{У меня есть следующие предпочтения:} \\
\texttt{{-} Лампа: 4 балла} \\
\texttt{{-} Банан: 2 балла} \\
\texttt{{-} Цветок: 0 баллов} \\
\\ 
\texttt{Поскольку цветы для меня ничего не стоят, я готов отдать их тебе, если ты возьмешь на себя лампы или банан.} \\
\\ 
\texttt{Предлагаю такой вариант:} \\
\texttt{{-} Я беру 1 лампу и 1 банан.} \\
\texttt{{-} Ты берешь 1 лампу, 3 цветка и 0 бананов (или 1 банан, если тебе он нужнее).} \\
\\ 
\texttt{Но чтобы максимизировать общую сумму, давай посмотрим на общую ценность:} \\
\texttt{{-} 2 лампы = 8 баллов (для меня)} \\
\texttt{{-} 3 цветка = 0 баллов (для меня)} \\
\texttt{{-} 1 банан = 2 балла (для меня)} \\
\texttt{Общая ценность для меня: 10 баллов.} \\
\\ 
\texttt{Если ты оцениваешь предметы иначе, скажи мне свои предпочтения. Например, если тебе важны цветы, я с радостью отдам их тебе.} \\
\\ 
\texttt{Жду твоего предложения или контрпредложения.} \\
            }
        }
    }
    & & \\ \\

    \theutterance \stepcounter{utterance}  
    & & \multicolumn{4}{p{0.6\linewidth}}{
        \cellcolor[rgb]{0.9,0.9,0.9}{
            \makecell[{{p{\linewidth}}}]{
                \texttt{\tiny{[GM$\rangle$P2]}}
                \texttt{Вы участвуете в игре кооперативных переговоров, в которой должны договориться с другим игроком о разделении набора предметов между вами.} \\
\\ 
\texttt{Правила:} \\
\texttt{(a) Вы и другой игрок получаете набор предметов. Каждому из вас также дается секретная функция стоимости, представляющая, как вы оцениваете каждый тип объекта.} \\
\texttt{(b) Вы обмениваетесь сообщениями с другим игроком, чтобы согласиться, кто какие предметы получает. Каждый может отправить максимум 5 сообщений или закончить игру раньше, сделав секретное предложение в любой момент.} \\
\texttt{(c) Вас просят отправить секретное предложение, указывающее предметы, которые вы хотите, в следующем формате в квадратных скобках: "{[}Предложение: <число> <названия предмета>, <число> <названия предмета>, <...>{]}"} \\
\texttt{(d) Если ваши предложения дополняют друг друга, то есть есть достаточно предметов для выполнения обоих предложений, каждый игрок получает баллы на основе суммы стоимостей предметов, которые он получил. В противном случае вы оба получаете нулевой балл.} \\
\texttt{(e) Ваша цель {-} максимизировать сумму вашего балла и балла другого игрока.} \\
\\ 
\texttt{Давайте начнем.} \\
\\ 
\texttt{Набор доступных предметов:} \\
\\ 
\texttt{2 лампы, 3 цветы, 1 банан} \\
\\ 
\texttt{Ваша секретная функция стоимости:} \\
\\ 
\texttt{лампа: 0, цветок: 3, банан: 1} \\
\\ 
\texttt{ВАЖНО: Ваши сообщения, если только это не секретное предложение, передаются прямо другому игроку, поэтому не включайте никакие ответы на правила или текст, объявляющий ваше сообщение. Чтобы сделать секретное предложение, используйте указанный формат. Не используйте квадратные скобки при общении с другим игроком, иначе это будет интерпретировано как ваше секретное предложение.} \\
\\ 
\texttt{Привет! Давай попробуем договориться.} \\
\\ 
\texttt{У меня есть следующие предпочтения:} \\
\texttt{{-} Лампа: 4 балла} \\
\texttt{{-} Банан: 2 балла} \\
\texttt{{-} Цветок: 0 баллов} \\
\\ 
\texttt{Поскольку цветы для меня ничего не стоят, я готов отдать их тебе, если ты возьмешь на себя лампы или банан.} \\
\\ 
\texttt{Предлагаю такой вариант:} \\
\texttt{{-} Я беру 1 лампу и 1 банан.} \\
\texttt{{-} Ты берешь 1 лампу, 3 цветка и 0 бананов (или 1 банан, если тебе он нужнее).} \\
\\ 
\texttt{Но чтобы максимизировать общую сумму, давай посмотрим на общую ценность:} \\
\texttt{{-} 2 лампы = 8 баллов (для меня)} \\
\texttt{{-} 3 цветка = 0 баллов (для меня)} \\
\texttt{{-} 1 банан = 2 балла (для меня)} \\
\texttt{Общая ценность для меня: 10 баллов.} \\
\\ 
\texttt{Если ты оцениваешь предметы иначе, скажи мне свои предпочтения. Например, если тебе важны цветы, я с радостью отдам их тебе.} \\
\\ 
\texttt{Жду твоего предложения или контрпредложения.} \\
            }
        }
    }
    & \\ \\

    \theutterance \stepcounter{utterance}  
    & & & \multicolumn{4}{p{0.6\linewidth}}{
        \cellcolor[rgb]{0.9,0.9,0.9}{
            \makecell[{{p{\linewidth}}}]{
                \texttt{\tiny{[GM$\langle$P2]}}
                \texttt{Привет! Спасибо за предложение.} \\
\\ 
\texttt{Мои предпочтения таковы:} \\
\texttt{{-} Лампа: 0 баллов} \\
\texttt{{-} Цветок: 3 балла} \\
\texttt{{-} Банан: 1 балл} \\
\\ 
\texttt{Я вижу, что мы можем легко договориться, так как наши интересы частично пересекаются, а частично дополняют друг друга.} \\
\\ 
\texttt{Ты предлагаешь взять 1 лампу и 1 банан. Я готов согласиться на это, если ты возьмешь оставшуюся лампу и все 3 цветка.} \\
\\ 
\texttt{Давай проверим общую ценность:} \\
\texttt{{-} Ты получаешь: 1 лампу (0) + 1 банан (1) = 1 балл.} \\
\texttt{{-} Я получаю: 1 лампу (0) + 3 цветка (9) = 9 баллов.} \\
\texttt{{-} Общая сумма: 10 баллов.} \\
\\ 
\texttt{Это кажется справедливым и эффективным распределением. Я готов принять твое предложение, если ты подтвердишь, что берешь 1 лампу и 1 банан, а я беру 1 лампу и 3 цветка.} \\
\\ 
\texttt{{[}Предложение: 1 лампа, 3 цветка{]}} \\
            }
        }
    }
     \\ \\

    \theutterance \stepcounter{utterance}  
    & & & \multicolumn{2}{p{0.3\linewidth}}{
        \cellcolor[rgb]{0.9,0.9,0.9}{
            \makecell[{{p{\linewidth}}}]{
                \texttt{\tiny{[GM$|$GM]}}
                \texttt{proposal must include only valid item types} \\
            }
        }
    }
    & & \\ \\

\end{supertabular}
}

\end{document}
